\subsection{From Interactive Uncertainty to Two Capabilities}
\label{sec:coordination-formulation}

Coordination with private information is naturally an interactive partially
observable decision problem.  For agent $i$, an I-POMDP is $\mathcal I_i=\langle IS_i,A,T_i,\Omega_i,O_i,R_i\rangle, IS_i=S\times\prod_{j\neq i}M_j$
where $S$ is the physical state, $M_j$ models agent $j$, $A$ is the joint
action space, and $T_i,O_i,R_i$ specify transition, observation, and
reward~\citep{gmytrasiewicz2005ipomdp}.  Uncertainty therefore covers not only
the world but also others' information, preferences, and policies.  This view
underlies opponent modelling and belief-space coordination
\citep{carmel1996opponent,huang2024hop,goel2026bayesian}; the observation--belief--action
links nevertheless remain fragile for LLM agents
\citep{sobotka2026strategic,jha2026modeling}. 
Let $b_t^i\in\Delta(IS_i)$ be agent $i$'s belief after public history $h_t$.
An observation updates it as $b_{t+1}^i=\tau_i(b_t^i,a_t^i,o_{t+1}^i)$, while
an action is selected by belief-conditioned planning, $a_t^{i,*}\in\arg\max_{a\in A_i}\mathbb E\!\left[\sum_{k=t}^{T}\gamma^{k-t}R_i(x_k,a_k)\mid b_t^i,a_t^i=a,\pi_{-i}\right]$
We call these computations \emph{partner inference} ($B$), which maintains the
part of $b_t^i$ predictive of others, and \emph{partner-conditioned planning}
($P$), which evaluates actions under that belief.  They form the loop
$B_t\!\rightarrow P_t\!\rightarrow a_t^i\!\rightarrow o_{t+1}^i
\!\rightarrow B_{t+1}$: belief changes the best interaction, while the chosen
interaction determines the next evidence.  This is a functional decomposition,
not a claim about separate neural modules.

\subsection{BENAC-I: Negotiation with Private Preferences}
\label{sec:benac-i}

To expose both computations while retaining exact supervision, we extend the
commitment-based negotiation game BENAC~\citep{benac2026benchmark}.  We call
the resulting incomplete-information variant \textsc{BENAC-I}.  An instance is
\begin{equation}
 \mathcal E=\langle\mathcal N,\mathbf K,\mathcal G,\Gamma,\delta,
 \rho,C_0,\mathcal L,\{\Theta_i\}_{i\in\mathcal N},\mu_0,\mathbf q\rangle .
 \label{eq:benac-i-game}
\end{equation}
Here $\mathcal N$ is the player set, $K_i$ counts $i$'s binary commitment
coordinates, and each public goal $g\in\mathcal G$ has requirements
$\Gamma_g\subseteq\{(i,k):k<K_i\}$ and mode
$\delta_g\in\{\mathrm{binary},\mathrm{linear}\}$.  The proposer schedule
$\rho$ and legal-action relation $\mathcal L$ govern play from the initially
empty, irreversible commitment state $C_0$.  Each player privately observes a
preference type $v_i\in\Theta_i\subseteq\{-1,0,+1\}^{|\mathcal G|}$, drawn
from common prior $\mu_0$; entries mean \textsc{Avoid}, \textsc{Neutral}, or
\textsc{Want}.  Vector $\mathbf q$ records remaining investigation budgets.
For commitment state $C$, goal satisfaction and terminal utility are
\begin{equation}
 S_g(C)=
 \begin{cases}
 \mathbf 1[\,c_{ik}=1,\ \forall(i,k)\in\Gamma_g\,],&\delta_g=\mathrm{binary},\\
 |\Gamma_g|^{-1}\sum_{(i,k)\in\Gamma_g}c_{ik},&\delta_g=\mathrm{linear},
 \end{cases}
 \qquad U_i(C;v_i)=\sum_g v_{ig}S_g(C).
 \label{eq:benac-i-utility}
\end{equation}
On a proposal turn, a player may \textsc{pass}, \textsc{offer} new commitments, or use
\textsc{Investigate}$(j,g)$ when $q_i>0$.  Acceptance irreversibly binds both
sides' proposed additions; rejection binds none; choosing from a two-offer
proposal binds only the selected option.  Investigation instead consumes the
turn and one unit of $q_i$, leaves $C$ unchanged, and privately reveals $v_{jg}$;
the target, but not the answer, is public.  Negotiation therefore supplies
indirect evidence and investigation supplies costly direct evidence within the
same action space.

\paragraph{Reference teacher and labels.}
The teacher is an information-limited policy profile, not an omniscient player.
At an information set $I_j=(h,v_j,r_j)$---public history, own preferences, and
private investigation results---it chooses
\begin{equation}
 \operatorname{supp}\pi_j^{\star}(\cdot\mid I_j)\subseteq
 \arg\max_{a\in\mathcal L(I_j)}
 \mathbb E_{v_{-j}\sim b_j(\cdot\mid I_j)}
       [U_j(C_T;v_j)\mid a,\pi_{-j}^{\star}],
 \label{eq:teacher-policy}
\end{equation}
breaking own-utility ties by total utility to others and mixing over residual
ties.  We enumerate the finite game tree, initialize legal actions uniformly,
and synchronously update information-set best responses until the whole
contingent profile is stable.  We retain it only after a unilateral-deviation
check; cycles and resource-limit failures yield no label.  This defines a
reproducible reference, not a unique or general-sum Nash equilibrium.

The same profile supplies both labels.  An observed action gives the exact
finite-support update $b_{t+1}^i(v_{-i})\propto b_t^i(v_{-i})
\pi_j^{\star}(u_t^j\mid I_t^j,v_j)$, treating $i$'s actions as interventions
and conditioning on any private investigation result; this is the $B$ target.
For $P$, exact continuation gives $Q_i^{\star}(I_i,a)$ for every legal action
and hence every optimal action.  Both labels thus share one declared behavioral
model and terminal objective.
